\section{Isocontour Extraction via Delaunay Triangulation}

\subsection{Introduction}
The extraction of isocontours from unstructured point clouds plays a crucial role in numerical simulation, visualization, and computational geometry. In the absence of a structured grid, Delaunay triangulation provides a rigorous framework to infer spatial connectivity among sample points. This section presents the theoretical basis for identifying isocontours in an \(N\)-dimensional domain, where each sample point is assigned a scalar potential. These values define a scalar field over the domain, enabling the detection of contour surfaces via interpolation.

\subsection{Point Cloud Structure and Scalar Field Definition}
Let \( X \in \mathbb{R}^{M \times (N+1)} \) be a matrix of \( M \) sample points. The first \( N \) columns of \( X \), denoted \( X_{:,1:N} \), specify the coordinates of each point in an \( N \)-dimensional Euclidean space. This space is assumed to have a Cartesian coordinate system with orthogonal axes, enabling the use of standard distance metrics. The last column, \( X_{:,N+1} \), assigns a scalar potential to each point and thereby defines a scalar field over the domain. This scalar field serves as the basis for extracting isocontours corresponding to a prescribed potential value \( \phi \in \mathbb{R} \).

\subsection{Delaunay Triangulation in \texorpdfstring{\( \mathbb{R}^N \)}{R\^N}}
From the base space \( X_{:,1:N} \), we compute the Delaunay triangulation \( \mathcal{T} \), which partitions the convex hull of the points into non-overlapping \(N\)-simplices. Each simplex is formed by \( N+1 \) affinely independent points and satisfies the empty circumsphere property: no other point lies inside the circumscribed sphere of any simplex. This ensures that the triangulation avoids sliver elements and provides a high-quality mesh suitable for interpolation.

\subsection{Isocontour Extraction via Edge-Wise Interpolation}
Each simplex generated by the triangulation defines a set of \( \binom{N+1}{2} \) edges, corresponding to all pairwise vertex combinations. For every such edge connecting vertices \( i \) and \( j \), we consider whether the isocontour intersects the edge by comparing the scalar potentials:
\begin{equation}
\phi_i = X(i, N+1), \quad \phi_j = X(j, N+1).
\end{equation}
An intersection is present if the scalar potential crosses the isovalue \( \phi \) along the edge, i.e.,
\begin{equation}
(\phi_i - \phi)(\phi_j - \phi) < 0.
\end{equation}
In that case, we compute the intersection point using linear interpolation.

\subsubsection*{Linear Interpolation in \texorpdfstring{\( \mathbb{R}^{N+1} \)}{R\^(N+1)}}
Assuming linear variation of the scalar field between the endpoints, the interpolation parameter \( \lambda \in [0,1] \) is given by:
\begin{equation}
\lambda = \frac{\phi - \phi_i}{\phi_j - \phi_i}.
\end{equation}
The location of the intersection point in the full \((N+1)\)-dimensional space is then computed as:
\begin{equation}
x_{\text{iso}} = (1 - \lambda) \cdot X(i,:) + \lambda \cdot X(j,:).
\end{equation}
By construction, the scalar potential component of \( x_{\text{iso}} \) satisfies \( x_{\text{iso}}(N+1) = \phi \), while the first \( N \) components represent the spatial location in the base space where the isocontour intersects the edge.
